\documentclass[12pt, norsk, a4paper]{proc}
\pagestyle{empty}

\usepackage[top=2cm, bottom=3cm, left=2cm, right=2cm]{geometry} % kan fjernes (bare brukt for å få litt bedre layout en gang)
\usepackage{setspace}
\usepackage{hyperref}
\usepackage{graphics}
\usepackage{cuted}
\usepackage{graphicx}
\usepackage{epstopdf}
\usepackage{wrapfig}
\usepackage[export]{adjustbox}
\usepackage{titlesec}


\label{cybernetisk-selskab}

\setlength{\parindent}{0pt} 
\onehalfspacing

\providecommand{\tightlist}{%
  \setlength{\itemsep}{0pt}\setlength{\parskip}{0pt}}
  
  
%----------------------------------------------------------------------------------------------
% TODO:
% ENDRE DISSE:
\newcommand{\genforsDato}{19. mai  kl. 17:00} % datoen og klokkeslettet genfors faller på
\newcommand{\genforsDag}{torsdag} % dagen genfors faller på ref. mandag, tirsdag osv.
\newcommand{\genforsDatoUkeFoer}{12. mai} % 7 dager før \genforDato
\newcommand{\generalforsamling}{Ordinær Generalforsamling} % 7 dager før \genforDato

\newcommand{\leder}{Gyda} % fornavn på leder
\newcommand{\nestlederFornavn}{Peter} % fornavn på deg selv (nestleder)
\newcommand{\nestlederFulltNavn}{Peter Hjelle Petersen-Øverleir} % fullt navn på den som lager innkallingen (nestleder)
\newcommand{\mogul}{Dag} %fornavn på kjellermogul
%----------------------------------------------------------------------------------------------


\title{Cybernetisk Selskab}

\begin{document}

   \begin{strip}
       \begin{center}
           \includegraphics[scale=1, center]{cyb-logo.eps} 
           \vspace*{\stretch{1.0}}
            \Large\textbf{Cybernetisk Selskab}\\
            \textit{inviterer deg til}\\
            \large\textbf{\generalforsamling} \\
            \emph{den} \\
            \large \genforsDato \\
           \vspace*{\stretch{2.0}}
       \end{center}
   \end{strip}

\textbf{Velkommen til \generalforsamling!}
Cybernetisk Selskab avholder \textbf{\genforsDag \space \genforsDato} generalforsamling i Escape. Alle med betalt medlemskap innen \genforsDatoUkeFoer \space har stemmerett.\\

Studenter ved IFI og medlemmer kan stille til valg. Det er en lang rekke verv
som skal fylles, fra 
% TODO:
% verv som skal fylles:
Kasserer 
til 
% TODO:
% verv som skal fylles:
DJ-sjef.
% ^^^^^^
Til
sist i innkallingen finner du en kort beskrivelse av alle vervene.\\

Vi fortsetter denne gangen også med oversikt og skjema for å sende inn vedtektsforslag. Disse kan henholdsvis finnes på \textbf{\href{http://forslag.cyb.no/}{forslag.cyb.no}} og \textbf{\href{http://forslagskjema.cyb.no/}{forslagskjema.cyb.no}}. Fristen for innsending av forslag er,  for \generalforsamling, 48 timer i forveien. \\

Om du lurer på noe, er det bare å ta kontakt med 
meg (\nestlederFornavn), \leder \space (leder), \mogul \space (Kjellermogul) eller noen andre i cyb. Vi er tilgjengelige på
discord, slack og epostene: \href{mailto:nestleder@cyb.no}{nestleder@cyb.no}, \href{mailto:leder@cyb.no}{leder@cyb.no} og
\href{mailto:kjellermogul@cyb.no}{kjellermogul@cyb.no}.

Vel møtt.\\\\
Med vennlig hilsen\\
\nestlederFulltNavn\\
Nestleder

% TODO:
% OPPDATER FORELØPIG DAGSORDEN:
\hypertarget{foreluxf8pig-dagsorden}{%
\subsection*{Foreløpig dagsorden}\label{foreluxf8pig-dagsorden}}

\begin{itemize}
\tightlist
\item
  Semesterberetning
\item
  Regnskap
\item
  Budsjett
\item
  Kontingentfastsettelse
\item
  Valg
\item
  Tildeling av utmerkelser
\end{itemize}
% ^^^^^^^^^^^^^^

% OPPDATER verv til valg:
\hypertarget{verv-til-valg}{%
\subsection*{Verv til valg}\label{verv-til-valg}}

\hypertarget{hovedstyret}{%
\subsubsection*{Hovedstyret}\label{hovedstyret}}

\emph{Hovedstyret sørger for at foreningen i sin helhet driftes på en
god og forsvarlig måte. Alle verv i hovedstyre har varighet på to
semestre.}

\begin{itemize}
\tightlist
\item
  Leder
\item
  Nestleder
 \item
  Kasserer
\item
  Internansvarlig
\end{itemize}

% TODO:
% OPPDATER VERVENE OG BESKRIVELSENE SOM STÅR HER
\hypertarget{om-vervene-i-hovedstyret}{%
\subsubsection*{Om vervene i hovedstyret}\label{om-vervene-i-hovedstyret}}

\textbf{Leder} er leder av hovedstyret og overhodet i foreningen, og har dermed det overordnede ansvaret for alt som foregår i foreningen.

\textbf{Nestleder} er leders høyre
hånd og stedfortreder, og skal kunne tre inn og hjelpe leder om leder skulle være
utilgjengelig. 


% \textbf{Kjellermogul} er leder av Kjellerstyret og daglig leder i
% Escape. Og har dermed det overordnede ansvaret for alt som foregår på
% Escape.

\textbf{Kasserer} er leder av økonomigruppa og har det overordnede ansvaret for økonomien i
foreningen. 

% \textbf{Arrangementssjef} Arrangementssjef er leder for arrangementsgruppa som organiserer og gjennomfører ulike arrangementer gjennom året. 

% \textbf{Rekrutteringsansvarlig} har som ansvar å få inn nye frivillige i
% CYB. Rekrutteringsansvarlig setter potensielle interne i kontakt med
% gruppeledere.

\textbf{Internansvarlig} har som hovedansvar å følge opp de interne i
CYB og hjelper til at man har en god tone innad foreningen.

% \textbf{Promoteringssjef} leder promoteringsgruppen som har ansvar for å
% utarbeide grafisk materiell til arrangementene i foreningen og skape
% blest rundt disse.

% \textbf{X-sjef} -------.....-----+++++ osv...



\hypertarget{kjellerstyret}{%
\subsubsection*{Kjellerstyret}\label{kjellerstyret}}

\emph{Kjellerstyret har ansvar for den daglige driften av Escape.
Vervene i kjellerstyret varer i ett semester, foruten økonomiansvarlig 
som varer i to semestre.}

\begin{itemize}
\tightlist
\item
  Kjellernestleder
\item
  Barsjef
\item
  Kafésjef
\item
  Økonomiansvarlig %Kanskje ikke oppe til valg?
\item
  Innkjøpsansvarlig
\item
  Teknisk sjef
\item
  Utlånsansvarlig
\item
  DJ-sjef
\item
  Arrangementskoordinator
\end{itemize}
% ^^^^^^^^^^^^^^^^^^^^^^^^^^^^^^^^^^^^


% TODO:
% OPPDATER VERVENE OG BESKRIVELSENE SOM STÅR HER
\hypertarget{om-vervene-i-kjellerstyret}{%
\subsubsection*{Om vervene i kjellerstyret}\label{om-vervene-i-kjellerstyret}}

\textbf{Kjellernestleder} er i praksis nestleder i kjellerstyret, og bidrar
med ulike oppgaver ved behov. Nestleder er også kjellermoguls høyre
hånd og stedfortreder, og skal kunne tre inn dersom kjellermogul skulle være
utilgjengelig.

\textbf{Barsjef} er nestleder i Kjellerstyret, og er derfor
ansvarsperson etter Kjellermogul. Hovedansvaret ligger i drift av bar.

\textbf{Kafésjef} er ansvarlig for å drive kaféen på dagtid i Escape og
har tett kontakt og oppfølgning av de frivillige i kaféen.

\textbf{Økonomiansvarlig} jobber tett med kasserer og har ansvaret for
den daglige økonomidriften av Escape.

\textbf{Innkjøpsansvarlig} har ansvar for å kjøpe inn det som trengs til
driften av Escape.

\textbf{Teknisk sjef} er en altmuligperson som sørger for at alt
teknisk utstyr i Escape fungerer i tilfredsstillende grad.

\textbf{Utlånsansvarlig} organiserer utlån av Escape, og pleier også å
være CYBs representant i Studentkjellernes personalforening.

\textbf{DJ-sjef} har ansvar for opplæring, rekruttering og mobilisering
av DJ-er til kveldsarrangementer, samt håndheve musikkprofil og skape
god stemning.

\textbf{Arrangementskoordinator} er bindeledd mellom bargruppa,
arrangementsgruppa og kjellerstyret for arrangementer med hovedvekt på
spritbardrift.

\href{https://wiki.cyb.no/display/AKTIV/Vervoversikt}{Mer informasjon om vervene ligger også på CYBs internwiki.}
\end{document}